\subsection*{Genome-resolved functional analysis}
\label{methods:mag_lifeonair}
{\emergencystretch=3em
Shotgun metagenomes from the same collection (551 read-bearing libraries, \texttt{EQ-\allowbreak SHOTGUN-\allowbreak READ-\allowbreak LIBS-\allowbreak v1.0.0}, sequenced across 554 paired-end runs on Illumina NovaSeq~6000 [221 runs] and Illumina NovaSeq~X [333 runs] instruments) were assembled with MEGAHIT and binned with SemiBin2; DAS\_Tool reconciled bins, CheckM2 assessed quality, and dRep grouped 12,434 retained bins at 95\% average nucleotide identity \cite{li2015megahit,pan2023semibin2,sieber2018dastool,chklovski2023checkm2,olm2017drep}. The resulting catalogue of 5,438 species-representative genomes (\texttt{EQ-\allowbreak MAG-\allowbreak REFINED-\allowbreak v1.0.0}) met completeness $\ge 50\%$ and contamination $\le 10\%$, comprising 126 high-quality and 5,311 strict medium-quality genomes ($<10\%$ contamination) under MIMAG standards plus one boundary representative (\texttt{48Dr2\_\_\allowbreak self\_\allowbreak SemiBin\_25}, $64.61\%$ completeness and $10.0\%$ contamination). In the foundational data release, GTDB-Tk classified the genomes against GTDB r232 (\texttt{gtdb\_\allowbreak r232\_\allowbreak classification}), and Bakta v1.12.0 and eggNOG-mapper v2.1.13 supplied the release gene annotations (\texttt{EQ-\allowbreak PROT-\allowbreak BAKTA-\allowbreak v1.0.0}, 19,186,792 proteins). CoverM v0.7.0 estimated genome abundances by competitive read mapping (\texttt{EQ-\allowbreak COVERM-\allowbreak REL-\allowbreak ABUND-\allowbreak v1.0.0}). After excluding one library with zero mapped reads (\texttt{F5Dr1}, $100\%$ unmapped; leaving 550 mapped libraries) and four out-of-core libraries from Sites~62 and 64 (\texttt{62PRr2}, \texttt{64Dr3}, \texttt{64PRr2} and \texttt{64Sr2}), the functional analysis retained 546 core-transect profiles across 59 sites (489 unique specimen keys; \texttt{EQ-\allowbreak ECOL-\allowbreak SHOTGUN-\allowbreak CORE-\allowbreak 546-\allowbreak v1.0.0}); exact specimen matching to the 1,186-sample amplicon cohort identified 442 retained libraries across 58 sites (\texttt{EQ-\allowbreak ECOL-\allowbreak MATCHED-\allowbreak 442-\allowbreak v1.0.0}). Supplementary Section~\ref{si-sup:mag_integration} defines cohort exclusions, genome-quality sensitivity and the competitive archaeal mapping checks (\texttt{EQ-\allowbreak COVERM-\allowbreak ARCHAEA-\allowbreak COMP-\allowbreak v1.0.0}).

We searched 20,451,986 representative-genome proteins called by Pyrodigal (\texttt{EQ-\allowbreak PROT-\allowbreak PYRODIGAL-\allowbreak REP-\allowbreak v1.0.0}) against 106 microTrait profiles at their published trusted cutoffs and evaluated 41 published rules \cite{karaoz2022microtrait}. A separate KO projection used the frozen analytical representative-genome KO matrix (\texttt{rep\_\allowbreak ko\_\allowbreak v2.tsv}, \texttt{EQ-\allowbreak KO-\allowbreak LEGACY-\allowbreak ANALYTICAL-\allowbreak v0}; IBEX \texttt{04\_\allowbreak functional/\allowbreak repko\_\allowbreak clean/\allowbreak results}) to assess dependence on annotation method; neither \texttt{EQ-\allowbreak PROT-\allowbreak PYRODIGAL-\allowbreak REP-\allowbreak v1.0.0} nor \texttt{rep\_\allowbreak ko\_\allowbreak v2.tsv} is asserted to derive from the 19,186,792-protein foundational Bakta release product (\texttt{EQ-\allowbreak PROT-\allowbreak BAKTA-\allowbreak v1.0.0}; [PENDING: \texttt{DATA-\allowbreak PROVENANCE-\allowbreak UNVERIFIED-\allowbreak PROTEIN-\allowbreak EQUIVALENCE}; verify upstream gene-caller and mapper provenance of the 20,451,986 Pyrodigal protein set and frozen \texttt{rep\_\allowbreak ko\_\allowbreak v2.tsv} matrix relative to the 19,186,792-protein Bakta release product]). Genome trait calls were weighted by abundance, with both the mapped-community and original-read denominators retained. For paired compartment contrasts, we averaged replicates within each site, visit and compartment and then averaged paired visit differences within sites. Confidence intervals resampled ten contiguous geographic blocks. An exact genus-level decomposition separated changes between shared genera, changes within shared genera and compartment-exclusive contributions. Whereas total trait contrasts and bacterial MAG Hill diversity are unaffected by sub-domain taxonomy labels, 1,898 full taxonomy strings (including 1,124 genus and 220 order labels) differ between the frozen analytical taxonomy snapshot (\texttt{EQ-\allowbreak TAX-\allowbreak LEGACY-\allowbreak ANALYTICAL-\allowbreak v0}, \texttt{genome\_\allowbreak quality.tsv}) used in the broad 41-trait matrix and canonical \texttt{gtdb\_\allowbreak r232\_\allowbreak classification} in \texttt{EQ-\allowbreak MAG-\allowbreak REFINED-\allowbreak v1.0.0}. For the three focal guilds (BetL, AhpCF and group 1l $[\mathrm{NiFe}]$-hydrogenase), the complete genus-turnover decomposition was rerun under canonical \texttt{gtdb\_\allowbreak r232\_\allowbreak classification} against the immutable legacy baseline using the identical abundance matrices, trait calls and paired visit sets, with 1,000 resamples across the ten contiguous spatial blocks (seed \texttt{20261108}) and individual-site resampling as a secondary sensitivity check, while retaining the frozen paired trait intervals from the original 1,000- and 2,000-draw analyses (Supplementary Table~\ref{si-tab:canonical_turnover_3guilds}). PICRUSt2 and MAG-derived KO profiles (\texttt{rep\_\allowbreak ko\_\allowbreak v2.tsv}) were normalised separately to their full KO totals and compared on the 442 exact specimen-matched libraries across 8,860 shared KOs, using the same paired design.
\par}

Functional responses and genome-community summaries were analysed with an adaptation of the spatial hierarchy above, using campaign, compartment and standardised log read depth as fixed effects, with and without dune field. The model retained site, site--visit and site--visit--compartment random intercepts and the two-dimensional Mat\'ern field. Responses were standardised before fitting a Student-$t$ likelihood with ten degrees of freedom; coefficients were transformed back to their analysis scale. One metagenome (\texttt{F45PRr3}) lacked read-depth metadata in the covariate table, leaving 545 profiles for these models (\texttt{EQ-ECOL-\allowbreak SHOTGUN-\allowbreak MODEL-\allowbreak 545-\allowbreak v1.0.0}) and 441 in the matched-cohort sensitivity (\texttt{EQ-ECOL-\allowbreak MATCHED-\allowbreak MODEL-\allowbreak 441-\allowbreak v1.0.0}). Sensitivities varied the spatial-range prior, likelihood and genome-quality restriction. These functional fits used the available spatial and sampling metadata; chemistry covariates were confined to the amplicon models. Numerical diagnostics and exclusions are reported in Supplementary Section~\ref{si-sup:mag_integration}.

\subsection*{Trace-gas oxidation genes and their carriers}

Hydrogenase large subunits were typed against HydDB, and CoxL proteins were classified using form-specific motifs and reference alignments \cite{sondergaard2016hyddb,cordero2019atmospheric}. We quantified groups 1l, 1h and 2a and form-I CoxL separately. Other [NiFe]-hydrogenases and other CoxL sequences supplied control groups. Ten million forward reads per library were sampled and truncated to 100 nt before competitive translated searches against the marker references and 14 universal single-copy ribosomal proteins. Length-normalised marker counts divided by median ribosomal-protein coverage defined genome equivalents. The primary analysis required 60\% amino-acid identity and detection of at least ten ribosomal markers. Excess compartment effects were the marker log$_2$ ratio minus the mean log$_2$ ratio of the two control sets.

Paired trace-gas contrasts used replicate medians within visits, site-level resampling and multiplicity adjustment within declared test families. Hierarchical spatial fits also analysed sample-level log$_2$(genome equivalents + 0.01) and excess-marker responses. Genus-level and order-adjusted analyses related carrier depletion to predicted maximal growth potential from gRodon2 and taxonomic order \cite{weissman2021grodon} under the frozen analytical taxonomy snapshot (\texttt{EQ-TAX-\allowbreak LEGACY-\allowbreak ANALYTICAL-\allowbreak v0}; [PENDING: \texttt{TAXONOMY-\allowbreak SNAPSHOT-\allowbreak SENSITIVITY}; confirm order- and genus-adjusted \texttt{gRodon2} regression using canonical assignments]). Supplementary Section~\ref{si-sup:loa_methods} gives the complete classification, read-assignment benchmarks, contamination controls, growth models and sensitivity analyses. These functional hypotheses were developed after inspecting the observational data.

\subsection*{Full microTrait annotation sensitivity}

We used the complete microTrait profiles of the shared 5,438-genome resource \cite{rubalkhali_datapaper,karaoz2022microtrait} to assess dependence of the two focal surface associations on annotation implementation. This workflow called genes with Prodigal v2.6.3 in metagenomic mode and applied the pinned microTrait profiles, coverage filters and logical rules. BetL required its transporter profile; AhpCF required an AhpC profile together with AhpF. We joined the resulting calls by exact canonical genome identifier to the same CoverM matrix and compared them with the earlier profile-based calls. For each annotation, we averaged replicate profiles within site, visit and compartment, calculated surface-minus-subsurface differences within visits, and averaged these differences within sites. We retained the same 68 visits at 36 sites, used 1,000 bootstrap resamples of ten contiguous transect blocks (seed 20261108), and reproduced the earlier estimates and intervals before evaluating the new calls. Mapped-community and original-read denominators were kept separate. This analysis is an annotation sensitivity check; the earlier INLA spatial-model estimates retain their original trait inputs. The complete rule matrices, gene-call provenance, missing-model flags and continuous predictions are documented in the data descriptor. The specialised group 1l hydrogenase classification and separate gRodon2 analyses retain their own inputs and methods.
